% !TEX root = ../paper.tex

\section{Additional Permit Evidence}
\label{sec:appendix_permits}

This appendix reports additional permit evidence that complements the descriptive discussion in Section~\ref{sec:data} and the event study in Section~\ref{sec:empirical_results_permits}.

\subsection{Balance Before Redistricting}
\label{subsec:permit_balance}

Table~\ref{tab:permit_itt_balance} compares blocks that changed wards in 2015 with nearby blocks that remained in the same ward.
Before the new map took effect, the two groups were broadly similar in permit activity and neighborhood characteristics.
The joint test across all listed characteristics has a $p$-value of 0.333.

\input{../tasks/permit_itt_balance/output/permit_itt_balance_500ft.tex}

\subsection{Permit Event-Study Estimates}
\label{subsec:permit_event_study_estimates}

Table~\ref{tab:permit_event_study_did} summarizes the average change in annual permit counts during the first six years under the new ward map.
The combined model treats reassignment toward greater stringency and greater leniency as changes of equal size in opposite directions.
A second model allows the two directions to have different effects.
Low-discretion permits show a more muted response.

\begin{table}[htbp]
    \centering
    \caption{Permit Effects with Combined and Separate Reassignment Directions}
    \label{tab:permit_event_study_did}
    \input{../tasks/run_event_study_permit/output/permit_event_study_appendix_500ft.tex}

    \footnotesize
    \textit{Notes:} Poisson regressions use census-block-year observations from 2010--2020 within 500~ft of a ward boundary.
    Stringency scores are estimated using permits filed from 2006 through 2014.
    The sample retains blocks whose 2014 origin- and destination-ward incumbents remained in office when the new map took effect.
    It also requires at least one high-discretion permit during event times $t=-5$ through $t=-1$.
    Blocks are classified as assigned toward greater stringency, assigned toward greater leniency, or unchanged under the new map.
    The combined model imposes equal-and-opposite effects for the two reassignment directions, with unchanged blocks as the reference group; a negative coefficient means that greater stringency reduces permits.
    The separate model includes one indicator for each reassignment direction.
    Its half-difference row puts the two-direction estimates on the same scale as the combined estimate.
    The final test asks whether the separately estimated effects are equal in size and opposite in sign.
    Each model includes block fixed effects and ward-pair $\times$ year fixed effects.
    The estimator omits blocks with no permits of the relevant type in any sample year and ward-pair-year groups with no such permits.
    Standard errors, in parentheses, are clustered by ward pair.
    $^{*}p<0.10$, $^{**}p<0.05$, $^{***}p<0.01$.
\end{table}

\subsection{Allowing the Two Reassignment Directions to Differ}
\label{sec:appendix_permit_unconstrained}

Panels B and C of Figure~\ref{fig:permit_event_study} allow moves toward greater stringency and greater leniency to have different effects.
The two paths need not be mirror images, but their difference provides a comparable measure of how permit activity changes with stringency.
Table~\ref{tab:permit_event_study_did} shows that this comparison is nearly identical to the combined estimate.

Panel A of Figure~\ref{fig:permit_event_study_checks} shows the corresponding path for low-discretion permits.
The response is more muted than for high-discretion permits.

\subsection{Including Wards with Alderman Turnover}
\label{sec:appendix_permit_assigned}

Panel B includes blocks from wards that elected new aldermen in 2015.
The reassignment classification still uses the two aldermen serving in 2014.
For blocks whose alderman subsequently changed, this measures the change implied by the map rather than the stringency the block necessarily experienced under the newly elected alderman.
The annual number of high-discretion permits is 8.3\% lower, a somewhat smaller estimate than in the main sample.

\begin{figure}[htbp]
    \centering
    \begin{minipage}[t]{0.48\textwidth}
        \centering
        \textbf{Panel A: Low-discretion permits}\par
        \includegraphics[width=\linewidth]{../tasks/run_event_study_permit/output/permit_event_study_low_discretion_nosigns_stable_signed_500ft.pdf}
    \end{minipage}\hfill
    \begin{minipage}[t]{0.48\textwidth}
        \centering
        \textbf{Panel B: Including wards with alderman turnover}\par
        \includegraphics[width=\linewidth]{../tasks/run_event_study_permit/output/permit_event_study_high_discretion_all_signed_500ft.pdf}
    \end{minipage}
    \caption{Additional Permit Event-Study Checks}
    \figurenotes{Panel A uses low-discretion permits excluding signs.
    Panel B uses high-discretion permits and includes wards with alderman turnover; assignment still uses the 2014 incumbents.
    Both panels combine the two reassignment directions using Equation~\eqref{eq:event_study_permits}; negative coefficients indicate fewer permits under assignment toward greater stringency.
    Points show effects in log points relative to 2014; subtitles report pooled 2015--2020 coefficients on the same scale.
    The pooled effects correspond to reductions of 1.7\% and 8.3\% in annual permits, respectively.
    Tests that the pre-2015 effects are jointly zero give $p=0.075$ in Panel A and $p=0.666$ in Panel B.
    Shaded areas are 95\% confidence intervals based on standard errors clustered by ward pair.
    * $p\leq0.10$, ** $p\leq0.05$, *** $p\leq0.01$.}
    \label{fig:permit_event_study_checks}
    \label{fig:permit_event_study_low_discretion}
    \label{fig:permit_event_study_assigned}
\end{figure}

\subsection{Permit Summary Statistics and Volume}
\label{app:permit_summary_stats}

Table~\ref{tab:permit_summary_stats} compares the high- and low-discretion permit groups and reports their ward-month processing-time and permit-volume correlations.
High-discretion permits are much slower than lower-discretion permits: the mean processing time is 57.3 days versus 26.2 days, and the median is 30 days versus 6 days.
They also display more cross-alderman variation in processing times.
Ward-months with longer high-discretion processing times have fewer high-discretion permits.
Low-discretion permits show a weak positive relationship instead.

\begin{table}[htbp]
    \centering
    \textbf{Panel A: Processing-time summary statistics}\par\smallskip
    \input{../tasks/permit_summary_stats/output/permit_processing_time_high_vs_low_alderman_summary.tex}

    \vspace{1em}
    \textbf{Panel B: Ward-month processing time and permit volume}\par\smallskip
    \input{../tasks/permit_summary_stats/output/permit_processing_time_volume_correlation.tex}
    \caption{Summary Statistics for High- and Low-Discretion Permits}
    \label{tab:permit_summary_stats}
    \label{tab:permit_volume_correlation}

    \footnotesize
    \textit{Notes:} The sample uses permits issued at least one day after application, with applications begun through December 2022.
    The permit groups are defined in Section~\ref{subsec:permit_pipeline}.
    The alderman-level interquartile range is computed across alderman-specific mean processing times.
    In Panel B, each observation represents one ward in one month. The raw correlation is the Pearson correlation between mean processing time and permit volume.
    The elasticity is the coefficient from a regression of mean log processing time on log permit volume with ward and year-month fixed effects; standard errors are clustered by ward.
    The All row pools high- and low-discretion permits. * $p<0.10$, ** $p<0.05$, *** $p<0.01$.
\end{table}
